\abstract{%
% Exponential advances in HPC hardware enables profound scientific and industrial innovation, but performance overhead from synchronization and error recovery has become increasingly challenging.
% ``Best-effort'' approaches that improve efficiency by relaxing guarantees of correctness and determinism have emerged as a promising remedy.
% Here, we test the performance and scalability of fully-asynchronous, best-effort communication on existing, commercially-available HPC hardware.
%
% A first set of experiments tested whether best-effort communication strategies can benefit performance compared to the traditional perfect communication model.
% At high CPU counts, best-effort communication improved both the number of computational steps executed per unit time and the solution quality achieved within a fixed-duration run window.
% Computation-heavy benchmark workloads yielded the strongest scaling efficiency, achieving at 64 processes 92\% the update-rate of single-process execution.
% We observed a relative speedup of up to $7.8\times$ under communication-heavy workloads.
%
% Under the best-effort model, characterizing the distribution of quality of service across processing components and over time is critical to understanding the actual computation being performed.
% Additionally, a complete picture of scalability under the best-effort model requires analysis of how such quality of service fares at scale.
% To answer these questions, we designed and measured a suite of quality of service metrics: simulation update period, message latency, message delivery failure rate, and message delivery coagulation.
% Under a lower communication-intensivity benchmark parameterization, we found that median values for all quality of service metrics were stable when scaling from 64 to 256 processes.
% Under maximal communication intensivity, we found only minor --- and, in most cases, nil --- degradation in median quality of service.
%
% In an additional set of experiments, we tested the effect of an apparently faulty compute node on performance and quality of service.
% Despite extreme quality of service degradation among that node and its clique, median performance and quality of service remained stable.
%
% We used the Conduit C++ library for best-effort communication to perform reported experiments.
% Development of this library stemmed from a practical need for a general-purpose, prepackaged framework for best-effort communication.
% We hope that free availability of this library, which includes built-in tools to measure quality of service metrics, can facilitate broader incorporation of best-effort communication into HPC applications.
%
Developments in high-performance computing (HPC) technology continue to drastically increase quantities of available processor cores, driven in compute-dense hardware accelerator devices.
% In the context of evolutionary biology, these technologies offer new opportunities for digital experiments exploring cross-scale biological phenomena — such as many-species eco-evolutionary dynamics and evolutionary transitions in individuality (e.g., multicellularity, eusociality).
In the context of digital evolution, this explosive growth offers immense opportunity to advance both hypothesis-driven explorations of multi-scale biological phenomena and application-driven optimization targeting rugged, expensive problem domains.
Because available processing power is vastly distributed, however, practical challenges in memory, bandwidth, fault-tolerance, and synchronization arise in harnessing it --- a problem epitomized by emerging AI/ML hardware accelerator platforms.
% (e.g., by Tenstorrent, Groq, Graphcore, SambaNova, and Cerebras), which rely on highly-distributed on-chip dataflow architectures comprising hundreds or thousands of cores.
For instance, to pack 880,000 processors onto a single chip, Cerebras' Wafer-Scale Engine (WSE) platform provisions only 42kB of on-device memory per processor, requires multi-hop routing for point-to-point communication, and supports data export only from the chip periphery.
To contend with such resource-constraint, compounded by inevietable hardware faults and irregularities at scale, partial relaxations of the traditional deterministic computing paradigms may become a worthwhile engineering trade-off; digital evolution is uniquely positioned to explore such possibilities and, indeed, has a rich history of doing so.
Such best-effort relaxations may focus around data collection, a paradigm we term ``inferential observability,'' or be extended to underlying core computation.
% In the latter case, to mitigate potential bias and improve reproducibility, we argue that instrumentation on real-time performance of simulation on hardware can help to understand the underlying structure of the computation actually being performed.
Here, we review two case studies of best-effort computing infrastructure developed for digital evolution projects --- one using traditional HPC nodes and the other using the Cerebras Wafer-Scale Engine --- and show how best-effort computing can improve scalability, and how suggest possible instrumentation to identify and mitigate potential biases.
% Evolution is uniquely amenable to disturbances --- including major extinctions --- and stochasticity, and out of necessity rich methods exist to overcome such challenges in allied fields of ecology and evolutionary biology.
% Further, in particular artificial life is meant to capture general dynamics, rather than to model anything specific, and evolutionary computation is meant to optimize rather than to reflect nature.
% Shifting away from a deterministic computing paradign, however, introduces new challenges in managing reproducibility, variance, and bias.
% These emerging platforms exploit highly-distributed on-chip dataflow architectures to achieve extraordinary computational throughput, at the cost of exacerbating strain on memory and bandwidth capacities. straining existing paradigms for high-performance computing. advent of new HPC accelerators
}
